\section{Techniques}\label{sec:techniques}

\begin{wrapfigure}{R}{0.3\textwidth}
    \centering
    \includegraphics[width=0.25\textwidth]{images/techniques}
\end{wrapfigure}

Although CI is a principle-based system which doesn't follow a strict collection of well-defined techniques, there are certain reoccurring movements which could be considered as techniques.
Yet, they should not be regarded as something to be followed literally, and as long as one stays within the framework of the CI principles, any adaption is judged as correct.

\subsection{Lifts}\label{subsec:lifts}

The most common, and most impressive, ``spice'' added to a dance are lifts.
Whenever one center is lifted above another person's, a lift can be easily executed by ``pouring the weight''.
There are different kinds of lifts; typical ones are:

\begin{itemize}
    \item \textbf{Hip-lift}: Usually ones butt, lower belly or side is placed on the lower back of the other person.
    \item \textbf{Shoulder-lift}: The highest level of a lift when flying on the shoulder of another person.
\end{itemize}

Lifting might lead to potentially dangerous situations, which require us to pay especially attention to the following safety rules:

\begin{itemize}
    \item Not grabbing is a general guideline, using less the hands and more the torso, yet with lifting this becomes especially important; thus no interlocking of the arms.
    \item Allows keep a hollow back (a.k.a. ``good gorilla'') to provide a stable support for your partner.
    \item Always keep your head above your ass, otherwise your partner will feel unsafe and slide down to the front.
\end{itemize}

Read more about this in the chapter ``\nameref{sec:safety}''.

\subsection{Spirals}\label{subsec:spirals}

To keep the pathway continuation, spirals are the perfect tool to doing so.

\textbf{Exercise:} Connect two points on the body, and visual and move along a connecting line between those by using spirals around that line.
Feel free to also play with different planes.

\subsection{Negative Space}\label{subsec:negative-space}

Dancing in the so-called \gls{negativespace} basically means entering, usually by reaching into it with a limb, the \gls{kinesphere} of another person without, yet engaging in touch.
This is a common way to non-verbally invite someone for contact, to dance together, at a jam for example.

\subsection{Others}\label{subsec:others}

\begin{itemize}
    \item \textbf{Body-surfing}: One person rolls on the floor, the other (usually) perpendicularly rolling over.
    \item \textbf{Counter-balance}: As also done in partner acrobatics, where the center is, for once, oppositional, instead of centers being shared.
\end{itemize}